\documentclass[12pt]{article}
\usepackage[utf8]{inputenc}
\usepackage[T1]{fontenc}
\usepackage{hyperref}
\title{Al-Sayyid Bedouin Sign Language Language Resources for K-12 Curriculum Development}
\author{Compiled by Claude Sonnet 5.5 (claude-sonnet-5-5) for Andr\'as Kornai}
\date{}
\begin{document}
\maketitle

\section{Introduction}
Al-Sayyid Bedouin Sign Language (ABSL) is a village sign language of a Bedouin community in the Negev, Israel, where hereditary deafness is common. It arose in the community over about three generations. The research below comes mainly from the University of Haifa Sign Language Research Lab. No school curriculum was found.

\section{Primary Educational Resources}
No primers, textbooks or school programmes for ABSL were found.

\section{Linguistic and Reference Materials}
\begin{itemize}
\item University of Haifa Sign Language Research Lab, ABSL pages. They describe research on morphology, prosody, sentence structure and the lexicon, and a 300-entry dictionary with an introduction in Arabic, Hebrew and English. The lab is led by Wendy Sandler and Irit Meir. \url{https://signlab.haifa.ac.il/?p=179}
\item Haifa lab, ABSL overview page. \url{https://signlab.haifa.ac.il/2013/01/15/absl-1/}
\item Endangered Languages Project entry, which cites Sandler, Meir, Padden and Aronoff (2005) on the emergence of the language, and Kisch, Zeshan and de Vos (2012). \url{https://www.endangeredlanguages.com/lang/6998?hl=en}
\item Wikipedia, Al-Sayyid Bedouin Sign Language. \url{https://en.wikipedia.org/wiki/Al-Sayyid_Bedouin_Sign_Language}
\end{itemize}

\section{Religious and Community Resources}
No religious or community-organisation resources were found.

\section{Digital and Audio Resources}
\begin{itemize}
\item Wikitongues page for syy. The page listed no videos or lexicons when the search was run. \url{https://wikitongues.org/languages/syy/}
\end{itemize}

\section{Curriculum Development Context}
ABSL is a shared sign language used by deaf people of all ages in its community. Documentation is research-oriented. The 300-entry dictionary is the only teaching-adjacent resource found. A curriculum would need to build on Haifa lab materials, with the community, and in sign video. Deaf signers should lead the work.
\end{document}